% p2-07-dsi

\section{P2 §7. Discrete Scale Invariance as a Mechanism}

7. Discrete Scale Invariance as a Mechanism

               Figure (fig-p2-ch7-dsi). DSI triptych --- Self-similar hierarchy / lambda = phi regime /
                                               Log-periodic signature.

  7.1 DSI in Hierarchical RG Systems [L1 --- Mathematical fact; L5 ---
  4D application]
  In standard RG, the group of scale transformations is continuous: mu $\rightarrow$ et mu for all t $\in$
  R. DSI restricts this to a discrete group generated by mu $\rightarrow$ lambda mu for a fixed
  lambda > 1 (Sornette, Phys. Rep. 297 (1998) 239). The effective action at scale mu then
  satisfies

  S[mu] = S[lambdan mu0], n $\in$ Z,

  with no constraint at intermediate scales. The observable consequence is log-periodic
  behavior in physical quantities:

  xi(mu) = mu-Delta [A + Bcos((2pi ln(mu/mu0))/(lnlambda) + phi0) + …]. 19

  The preferred ratio lambda characterizes the DSI structure and is, in general, a free
  parameter fixed by the underlying geometry.

  7.2 When Can lambda = $\varphi$? [L2 --- quasicrystals; L5 --- 4D QFT]
  [L2 --- Established in quasicrystal context.] The Fibonacci substitution system has
  substitution matrix

  MFib = 1 \& 1 1 \& 0 20

  with Perron-Frobenius eigenvalue $\varphi$. Quasicrystals governed by Fibonacci tiling inherit
  a DSI ratio of $\varphi$ from this eigenvalue as an algebraic fact --- not a fitted parameter
  (Shechtman et al., Phys. Rev. Lett. 53 (1984) 1951; Levine and Steinhardt, Phys. Rev.
  Lett. 53 (1984) 2477).

  [L3 --- Model level in Toda context.] As established in Section 4.3, the mass spectra of
  Toda theories with H3 (icosahedral) symmetry contain $\varphi$ factors (Fring and Korff, Nucl.
  Phys. B 729 (2005) 361). If an RG flow maps between scales separated by the H3
  dilation factor, the DSI ratio inherits $\varphi$.

  [L5 --- Conjectural step.] The conjecture is that there exists a class of 4D field theories
  whose RG group is restricted, by some UV boundary condition or by the topology of the
  coupling-constant manifold, to discrete dilations with ratio $\varphi$. No derivation of this has
  been given. The identification of such a restriction with any known physical principle
  (compactification, anomaly matching, holographic duality) is an open problem
  constituting Work Package 3. The formal failure criterion is given in Section 4.4.

  7.3 Log-Periodic Signatures in Running Couplings [L5 --- Template
  only; no prediction]
  If a $\varphi$-DSI mechanism operates, it would produce log-periodic corrections to the
  running of gauge couplings. The corrected coupling at scale mu would read

  alphai-1(mu) = alphai-1(mu0) + (bi)/(2pi)ln(mu)/(mu0) + Ai cos((2pi ln(mu/mu0))/(ln$\varphi$)
  + deltai) + O(alphai), 21

  where the first two terms are the standard one-loop result (Wilson and Kogut, Phys.
  Rep. 12 (1974) 75) and the cosine term is the DSI correction. The period in log-scale is

  ln$\varphi$ approx 0.4812, corresponding to an energy ratio of $\varphi$ approx 1.618 between
  successive nodes of the oscillation.

  This is not a prediction. The amplitudes Ai are free parameters in this framework.
  Expression (21) is a template for what a $\varphi$-DSI signal would look like, not a derivation.

  Current observational bound. The absence of observed log-periodic oscillations in LEP,
  LHC, and Tevatron coupling measurements sets an upper bound Ai $\leq$sssim few $\times$ 10-3
  under the assumption that the DSI period is ln$\varphi$. This bound must be computed
  precisely as part of Work Package 4. Future lepton colliders (ILC, CLIC, FCC-ee)
  projecting per-mille coupling precision would extend the bound by roughly an order of
  magnitude and potentially constitute the definitive test of the $\varphi$-DSI hypothesis at
  accessible energy scales.

\subsection{References}

- Coldea, R. et al. "Quantum Criticality in an Ising Chain: Experimental Evidence for Emergent E8 Symmetry." \textit{Science} 327 (2010) 177. DOI \href{https://doi.org/10.1126/science.1180085}{10.1126/science.1180085}.
- Wu, J. et al. "E8 Symmetry in CoN$b_{2}$$O_{6}$ at Higher Resolution." \textit{Phys. Rev. B} 108, L060403 (2023). DOI \href{https://doi.org/10.1103/PhysRevB.108.L060403}{10.1103/PhysRevB.108.L060403}.
- Fring, A.; Korff, C. "Affine Toda field theories related to Coxeter groups of non-crystallographic type." \textit{Nucl. Phys. B} 729 (2005) 361. \href{https://arxiv.org/abs/hep-th/0506226}{arXiv:hep-th/0506226}. DOI \href{https://doi.org/10.1016/j.nuclphysb.2005.08.044}{10.1016/j.nuclphysb.2005.08.044}.
- Pellis, S. "The Fine-Structure Constant and the Golden Ratio." \href{https://vixra.org/pdf/2110.0117v4.pdf}{viXra 2110.0117 v4}.
